% ==============================================================================
% PRESENTACION - ACTIVIDAD 6 / FILOSOFIA DEL DERECHO / UnADM
% Tema: Caso Ayotzinapa, justicia, fines del Derecho y derechos humanos.
%
% Base de contenido:
% - reporte-filosofia-del-derecho-Actividad-6.tex
% - Planificacion de Actividades S8 - Filosofia del Derecho.pdf
%
% Compilar desde la raiz del repositorio:
% .\scripts\latexmk-build.ps1 .\UnADM\filosofia-del-derecho\presentacion-filosofia-del-derecho-Actividad-6.tex
% ==============================================================================

\documentclass[
  spanish,
  aspectratio=169,
  xcolor={dvipsnames,table}
]{beamer}

\geometry{paperwidth=19.2cm,paperheight=10.8cm}

\usepackage[utf8]{inputenc}
\usepackage[T1]{fontenc}
\usepackage[spanish,es-tabla]{babel}
\usepackage[scaled=.96]{helvet}
\renewcommand{\familydefault}{\sfdefault}
\usepackage{array}
\usepackage{booktabs}
\usepackage{graphicx}
\usepackage{ragged2e}
\usepackage{tabularx}
\usepackage{tikz}
\usepackage{hyperref}
\usetikzlibrary{arrows.meta,calc,fit,matrix,positioning,shapes.geometric,shadows.blur}

\newcolumntype{Y}{>{\RaggedRight\arraybackslash}X}

% ------------------------------------------------------------------------------
% DATOS DE LA ACTIVIDAD
% ------------------------------------------------------------------------------
\newcommand{\studentname}{Martin Jonathan de la Cruz}
\newcommand{\studentshort}{M. J. de la Cruz}
\newcommand{\studentid}{ES2611202040}
\newcommand{\studentemail}{martin.delacruz@unadmexico.mx}
\newcommand{\universityname}{Universidad Abierta y a Distancia de Mexico}
\newcommand{\facultyname}{Licenciatura en Derecho}
\newcommand{\coursename}{Filosofia del Derecho}
\newcommand{\coursecode}{030}
\newcommand{\activitytitle}{Analisis filosofico-juridico del caso Ayotzinapa}
\newcommand{\activitysubtitle}{Actividad 6 - Video academico}
\newcommand{\activityweek}{Semana 8}
\newcommand{\teachingfigure}{Mtra. Reyna Patricia Gonzalez Rodriguez}
\newcommand{\deliverydate}{07 de junio de 2026}
\newcommand{\locationname}{Roma Norte, Ciudad de Mexico}
\newcommand{\departmentlogo}{img/departamentos/UnADM.pdf}
\newcommand{\campusurl}{https://www.unadmexico.mx/}

% ------------------------------------------------------------------------------
% PALETA VISUAL
% ------------------------------------------------------------------------------
\definecolor{unadmGreenDark}{HTML}{174A3A}
\definecolor{unadmGreen}{HTML}{5F8F3A}
\definecolor{unadmGold}{HTML}{B88A2A}
\definecolor{unadmPaper}{HTML}{F6F7F2}
\definecolor{unadmInk}{HTML}{1F2A24}
\definecolor{unadmMuted}{HTML}{5D6B66}
\definecolor{unadmSoft}{HTML}{E9EEE6}

\mode<presentation>{
  \usetheme{default}
  \usefonttheme{professionalfonts}
  \setbeamertemplate{navigation symbols}{}
  \setbeamertemplate{blocks}[rounded][shadow=false]
  \setbeamercovered{invisible}
}
\setbeamersize{text margin left=0.55cm,text margin right=0.55cm}

\hypersetup{
  colorlinks=true,
  urlcolor=unadmGreenDark,
  linkcolor=unadmGreenDark
}
\Urlmuskip=0mu plus 1mu\relax

\setbeamercolor{background canvas}{bg=white}
\setbeamercolor{normal text}{fg=unadmInk,bg=white}
\setbeamercolor{structure}{fg=unadmGreenDark}
\setbeamercolor{frametitle}{fg=white,bg=unadmGreenDark}
\setbeamercolor{title}{fg=white,bg=unadmGreenDark}
\setbeamercolor{subtitle}{fg=white}
\setbeamercolor{block title}{fg=white,bg=unadmGreen}
\setbeamercolor{block body}{fg=unadmInk,bg=unadmPaper}
\setbeamercolor{alerted text}{fg=unadmGold}
\setbeamerfont{title}{size=\LARGE,series=\bfseries}
\setbeamerfont{frametitle}{size=\large,series=\bfseries}
\setbeamerfont{block title}{series=\bfseries}
\setbeamertemplate{itemize item}{\textcolor{unadmGreen}{\large$\blacktriangleright$}}
\setbeamertemplate{itemize subitem}{\textcolor{unadmGold}{\scriptsize$\blacksquare$}}
\setbeamertemplate{enumerate item}{\textcolor{unadmGreenDark}{\insertenumlabel.}}

\setbeamertemplate{frametitle}{
  \nointerlineskip
  \begin{beamercolorbox}[wd=\textwidth,ht=1.02cm,dp=0.22cm,leftskip=0.35cm,rightskip=0.25cm]{frametitle}
    \begin{minipage}[c]{0.76\textwidth}
      \insertframetitle
    \end{minipage}
    \hfill
    \raisebox{-0.1cm}{\includegraphics[height=0.60cm]{\departmentlogo}}
  \end{beamercolorbox}
  \vspace{-0.04cm}
  {\color{unadmGold}\rule{\textwidth}{1.3pt}}
}

\setbeamertemplate{footline}{
  \leavevmode%
  \hbox{%
    \begin{beamercolorbox}[wd=.34\paperwidth,ht=2.7ex,dp=1ex,leftskip=1em]{author in head/foot}%
      \color{white}\studentshort
    \end{beamercolorbox}%
    \begin{beamercolorbox}[wd=.40\paperwidth,ht=2.7ex,dp=1ex,center]{title in head/foot}%
      \color{white}\coursename
    \end{beamercolorbox}%
    \begin{beamercolorbox}[wd=.26\paperwidth,ht=2.7ex,dp=1ex,rightskip=1em plus 1fil]{date in head/foot}%
      \color{white}\coursecode\hfill\insertframenumber/\inserttotalframenumber
    \end{beamercolorbox}%
  }%
  \vskip0pt%
}
\setbeamercolor{author in head/foot}{bg=unadmGreenDark}
\setbeamercolor{title in head/foot}{bg=unadmGreen}
\setbeamercolor{date in head/foot}{bg=unadmGreenDark}

\newcommand{\accentbar}{%
  \begin{tikzpicture}[remember picture,overlay]
    \path[use as bounding box] (0,0) rectangle (0,0);
    \fill[unadmGreen] (current page.south west) rectangle ([xshift=0.56\paperwidth,yshift=0.08cm]current page.south west);
    \fill[unadmGold] ([xshift=0.56\paperwidth]current page.south west) rectangle ([xshift=\paperwidth,yshift=0.08cm]current page.south west);
  \end{tikzpicture}
}

\newcommand{\sectionframe}[1]{%
  \begin{frame}[plain]
    \begin{tikzpicture}[remember picture,overlay]
      \path[use as bounding box] (0,0) rectangle (0,0);
      \fill[unadmGreenDark] (current page.south west) rectangle (current page.north east);
      \node[opacity=0.09,anchor=east] at ([xshift=-0.45cm]current page.east) {\includegraphics[width=0.43\paperwidth]{\departmentlogo}};
      \fill[unadmGreen] ([xshift=1.05cm,yshift=2.25cm]current page.south west) rectangle ([xshift=1.25cm,yshift=6.40cm]current page.south west);
      \fill[unadmGold] ([xshift=1.48cm,yshift=2.25cm]current page.south west) rectangle ([xshift=1.68cm,yshift=5.25cm]current page.south west);
      \node[
        anchor=west,
        align=left,
        text=white,
        text width=0.58\paperwidth,
        font=\fontsize{26}{31}\selectfont\bfseries
      ] at ([xshift=2.08cm,yshift=0.55cm]current page.west) {#1};
      \node[
        anchor=west,
        align=left,
        text=white!78,
        font=\large
      ] at ([xshift=2.08cm,yshift=-0.55cm]current page.west) {\activitysubtitle};
    \end{tikzpicture}
  \end{frame}
}

\newcommand{\unadmsection}[1]{%
  \section{#1}
  \sectionframe{#1}
}

\newcommand{\pill}[2]{%
  \begin{tikzpicture}[baseline=(n.base)]
    \node[
      rounded corners=2pt,
      fill=#1,
      text=white,
      inner xsep=5pt,
      inner ysep=2pt,
      font=\scriptsize\bfseries
    ] (n) {#2};
  \end{tikzpicture}%
}

\newcommand{\sourcecite}[1]{%
  \vspace{0.08cm}
  {\tiny\color{unadmMuted}\RaggedRight Cita: #1\par}
}

% En diapositivas conviene usar cita APA breve dentro del texto,
% por ejemplo Autor (anio) o (Autor, anio), y reservar la referencia
% completa para las diapositivas finales.

\title[\coursecode]{\activitytitle}
\subtitle{\activitysubtitle}
\author[\studentshort]{\texorpdfstring{\studentname\\\footnotesize Matricula: \studentid\\\href{mailto:\studentemail}{\studentemail}}{\studentname, Matricula: \studentid}}
\institute[UnADM]{\universityname\\\facultyname\\\coursecode\ \textperiodcentered\ \coursename}
\date[\activityweek]{\locationname\\\activityweek\\\deliverydate}

\begin{document}

% Portada institucional
\begin{frame}[plain]
  \begin{tikzpicture}[remember picture,overlay]
    \path[use as bounding box] (0,0) rectangle (0,0);
    \fill[unadmGreenDark,opacity=0.96] (current page.south west) rectangle ([xshift=0.58\paperwidth]current page.north west);
    \fill[unadmGreen,opacity=0.92] ([xshift=0.58\paperwidth]current page.south west) rectangle (current page.north east);
    \node[anchor=center,opacity=0.10] at (current page.center) {\includegraphics[width=0.72\paperwidth]{\departmentlogo}};
    \node[anchor=north west] at ([xshift=0.58cm,yshift=-0.46cm]current page.north west) {\includegraphics[width=2.55cm]{\departmentlogo}};
    \fill[unadmGold] ([xshift=0.60\paperwidth,yshift=1.18cm]current page.south west) rectangle ([xshift=0.93\paperwidth,yshift=1.30cm]current page.south west);
    \fill[white,opacity=0.15] ([xshift=0.60\paperwidth,yshift=0.92cm]current page.south west) rectangle ([xshift=0.84\paperwidth,yshift=1.02cm]current page.south west);
  \end{tikzpicture}

  \vspace*{1.78cm}
  \begin{minipage}{0.55\paperwidth}
    {\color{white}\fontsize{24}{29}\selectfont\bfseries\raggedright \activitytitle\par}
    \vspace{0.22cm}
    {\color{white!86}\large \activitysubtitle}\\[0.32cm]
    {\color{unadmGold}\rule{0.82\linewidth}{1.3pt}}\\[0.42cm]
    {\color{white}\studentname}\\
    {\color{white!82}\footnotesize Matricula: \studentid}
  \end{minipage}
  \hfill
  \begin{minipage}{0.34\paperwidth}
    \raggedleft
    {\color{white}\Large\bfseries UnADM}\\[0.18cm]
    {\color{white!84}\footnotesize \universityname}\\[0.45cm]
    {\color{white}\small \facultyname}\\
    {\color{white!78}\small \coursename}\\[0.45cm]
    {\color{white!74}\footnotesize \coursecode\ \textperiodcentered\ \activityweek}\\
    {\color{white!74}\footnotesize \teachingfigure}\\
    {\color{white!74}\footnotesize \locationname}\\
    {\color{white!74}\footnotesize \deliverydate}
  \end{minipage}
\end{frame}

\begin{frame}{Ruta del video academico}
  \begin{columns}[T]
    \begin{column}{0.48\textwidth}
      \begin{block}{Objetivo de la actividad}
        Analizar la relacion entre justicia, fines del Derecho y derechos
        humanos mediante el caso Ayotzinapa, para valorar como el Estado de
        Derecho responde ante un problema juridico contemporaneo en Mexico.
      \end{block}
      \begin{alertblock}{Tesis}
        Ayotzinapa no prueba solo una falla penal: prueba si el derecho opera
        como limite moral del poder y garantia real de dignidad.
      \end{alertblock}
    \end{column}
    \begin{column}{0.48\textwidth}
      \begin{block}{Organizacion tematica}
        \begin{itemize}
          \item Caso y problema central.
          \item Justicia, equidad y normas aplicables.
          \item Seguridad juridica y derechos humanos.
          \item Propuestas, conclusion y fuentes.
        \end{itemize}
      \end{block}
    \end{column}
  \end{columns}
  \sourcecite{(Universidad Abierta y a Distancia de Mexico [UnADM], 2026).}
  \accentbar
\end{frame}

\begin{frame}{Contenido}
  \tableofcontents
\end{frame}

\unadmsection{Caso y problema}

\begin{frame}{Caso elegido: Ayotzinapa}
  \begin{columns}[T]
    \begin{column}{0.46\textwidth}
      \begin{block}{Hecho base}
        Entre la noche del 26 y la madrugada del 27 de septiembre de 2014,
        43 estudiantes de la Escuela Normal Rural ``Raul Isidro Burgos''
        desaparecieron en Iguala, Guerrero.
      \end{block}
      \begin{block}{Por que es relevante}
        Como documenta la CNDH (2024), el caso concentra privacion de la
        libertad, afectacion a integridad, dolor prolongado de victimas
        indirectas y fallas de busqueda, investigacion, verdad y justicia.
      \end{block}
    \end{column}
    \begin{column}{0.50\textwidth}
      \begin{alertblock}{Problema filosofico-juridico}
        Que ocurre cuando el propio aparato estatal, por accion, omision o
        encubrimiento, deja de operar como garantia y se convierte en fuente
        de vulneracion.
      \end{alertblock}
      \begin{block}{Lectura del reporte}
        El caso debe analizarse como cruce entre validez juridica,
        responsabilidad etica y legitimidad institucional.
      \end{block}
    \end{column}
  \end{columns}
  \sourcecite{(Comision Nacional de los Derechos Humanos [CNDH], 2024; Camara de Diputados, 2026b).}
  \accentbar
\end{frame}

\begin{frame}{Del hecho al problema de Derecho}
  \begin{center}
    \resizebox{0.94\linewidth}{!}{%
    \begin{tikzpicture}[
      node distance=0.65cm,
      box/.style={
        rounded corners=4pt,
        draw=unadmGreen!35,
        fill=unadmPaper,
        text width=3.15cm,
        minimum height=1.05cm,
        align=center,
        font=\small
      },
      arrow/.style={-{Stealth[length=2.4mm]},very thick,draw=unadmGold}
    ]
      \node[box] (h) {\textbf{Hecho}\\ desaparicion y fallas de busqueda};
      \node[box,right=of h] (d) {\textbf{Derechos}\\ libertad, integridad, verdad y justicia};
      \node[box,right=of d] (e) {\textbf{Estado}\\ deber de prevenir, buscar e investigar};
      \node[box,right=of e] (p) {\textbf{Pregunta}\\ cuando la legalidad pierde legitimidad};
      \draw[arrow] (h) -- (d);
      \draw[arrow] (d) -- (e);
      \draw[arrow] (e) -- (p);
    \end{tikzpicture}
    }
  \end{center}
  \begin{alertblock}{Idea central}
    La Filosofia del Derecho no se limita a preguntar que normas existen:
    tambien revisa si esas normas justifican, orientan y limitan realmente al poder.
  \end{alertblock}
  {\small\color{unadmInk}Desde la lectura de Ruiz Rodriguez (2009) y Rojas-Gonzalez (2018), el caso obliga a pasar del hecho al criterio: no basta con nombrar la norma; hay que preguntarse si el sistema realmente limita al poder.}
  \sourcecite{(Ruiz Rodriguez, 2009; Rojas-Gonzalez, 2018; CNDH, 2024).}
  \accentbar
\end{frame}

\unadmsection{Justicia y fines del Derecho}

\begin{frame}{Justicia en la Filosofia del Derecho}
  \begin{columns}[T]
    \begin{column}{0.32\textwidth}
      \begin{block}{Como valor}
        Siguiendo a Ruiz Rodriguez (2009), la justicia permite evaluar si una
        norma o decision protege a la persona y no solo si fue emitida
        formalmente.
      \end{block}
    \end{column}
    \begin{column}{0.32\textwidth}
      \begin{block}{Como fin}
        El Derecho debe ordenar la convivencia, limitar el abuso y producir
        respuestas institucionales frente al dano.
      \end{block}
    \end{column}
    \begin{column}{0.32\textwidth}
      \begin{block}{En Ayotzinapa}
        La justicia exige verdad, busqueda, investigacion, reparacion y
        garantias de no repeticion.
      \end{block}
    \end{column}
  \end{columns}
  \vspace{0.25cm}
  \begin{alertblock}{Lectura critica}
    Una legalidad que no protege dignidad puede conservar forma juridica, pero
    pierde fuerza como razon publica de obediencia.
  \end{alertblock}
  \sourcecite{(Ruiz Rodriguez, 2009; Rojas-Gonzalez, 2018; Laun, 2021).}
  \accentbar
\end{frame}

\begin{frame}{Equidad, normas y solucion del conflicto}
  \small
  \begin{tabularx}{\linewidth}{@{}p{0.22\linewidth}Y Y@{}}
    \toprule
    \textbf{Criterio} & \textbf{Funcion en el caso} & \textbf{Salida esperada} \\
    \midrule
    Justicia & Reconocer dano, responsabilidad y deber de respuesta. & Verdad, sancion, reparacion y no repeticion. \\
    Equidad & Tratar de forma reforzada a victimas directas e indirectas. & Atencion diferenciada, informacion y acompanamiento. \\
    Norma mexicana & Activar Constitucion, Ley General de Victimas y Ley de Desaparicion. & Busqueda, investigacion, reparacion integral. \\
    Orden internacional & Leer el caso desde universalidad, debida diligencia y proteccion de derechos humanos. & Estandares de prevencion, verdad y reparacion. \\
    \bottomrule
  \end{tabularx}
  \vspace{0.25cm}
  \begin{alertblock}{Postura}
    Resolver el conflicto exige unir valor y procedimiento: justicia sin norma
    queda como deseo; norma sin justicia se vuelve formalismo.
  \end{alertblock}
  \sourcecite{(Camara de Diputados, 2024, 2026a, 2026b).}
  \accentbar
\end{frame}

\begin{frame}{Justicia distributiva y correctiva}
  \begin{columns}[T]
    \begin{column}{0.48\textwidth}
      \begin{block}{Justicia distributiva}
        En el caso, exige distribuir recursos, cargas y atencion publica de
        modo proporcional al dano: busqueda inmediata, investigacion seria,
        apoyo a familias y prioridad institucional.
      \end{block}
      \begin{alertblock}{Riesgo}
        Si el Estado distribuye indiferencia o burocracia, reproduce desigualdad
        frente a las victimas.
      \end{alertblock}
    \end{column}
    \begin{column}{0.48\textwidth}
      \begin{block}{Justicia correctiva}
        Busca responder al dano ya producido: esclarecer hechos, sancionar
        responsables, reparar integralmente y crear garantias de no repeticion.
      \end{block}
      \begin{alertblock}{Riesgo}
        Si no hay correccion, el dano se vuelve impunidad y el derecho pierde
        credibilidad como mecanismo de justicia.
      \end{alertblock}
    \end{column}
  \end{columns}
  \sourcecite{(Ruiz Rodriguez, 2009; Rojas-Gonzalez, 2018).}
  \accentbar
\end{frame}

\begin{frame}{Fines del Derecho puestos a prueba}
  \begin{center}
    \resizebox{0.88\linewidth}{!}{%
    \begin{tikzpicture}[
      every node/.style={font=\small},
      center/.style={ellipse,fill=unadmGreenDark,text=white,minimum width=3.3cm,minimum height=1.0cm,align=center},
      goal/.style={rounded corners=4pt,fill=unadmPaper,draw=unadmGreen!45,text width=3.2cm,minimum height=0.85cm,align=center},
      link/.style={-{Stealth[length=2.2mm]},draw=unadmGold,very thick}
    ]
      \node[center] (c) {Fines del Derecho};
      \node[goal,above left=1.0cm and 2.0cm of c] (j) {Justicia y equidad};
      \node[goal,above right=1.0cm and 2.0cm of c] (s) {Seguridad juridica};
      \node[goal,below left=1.0cm and 2.0cm of c] (b) {Bien comun y utilidad publica};
      \node[goal,below right=1.0cm and 2.0cm of c] (d) {Dignidad y derechos humanos};
      \draw[link] (c) -- (j);
      \draw[link] (c) -- (s);
      \draw[link] (c) -- (b);
      \draw[link] (c) -- (d);
    \end{tikzpicture}
    }
  \end{center}
  \begin{alertblock}{Aplicacion al caso}
    Ayotzinapa muestra que los fines del Derecho no son abstractos: como
    sostienen Ruiz Rodriguez (2009) y Rojas-Gonzalez (2018), se miden en la
    capacidad de proteger personas, controlar autoridades y reparar danos.
  \end{alertblock}
  \sourcecite{(Ruiz Rodriguez, 2009; Rojas-Gonzalez, 2018).}
  \accentbar
\end{frame}

\unadmsection{Normas y derechos}

\begin{frame}{Marco juridico relacionado}
  \small
  \begin{tabularx}{\linewidth}{@{}p{0.23\linewidth}Y Y@{}}
    \toprule
    \textbf{Fuente} & \textbf{Contenido relevante} & \textbf{Funcion en el caso} \\
    \midrule
    Constitucion mexicana & Articulos 1, 14, 16, 17 y 20. & Derechos humanos, seguridad juridica, debido proceso, acceso a justicia y derechos de victimas. \\
    Ley General de Victimas & Ayuda, asistencia, verdad, reparacion integral y garantias de no repeticion. & Traduce el dano en obligaciones concretas de atencion y reparacion. \\
    Ley General en Materia de Desaparicion & Busqueda, investigacion, coordinacion y respuesta institucional. & Convierte la desaparicion en deber reforzado del Estado. \\
    Estandares internacionales & Universalidad, debida diligencia, verdad, justicia y reparacion. & Evitan que el caso sea tratado como problema local aislado. \\
    \bottomrule
  \end{tabularx}
  \sourcecite{(Camara de Diputados, 2024, 2026a, 2026b; Suprema Corte de Justicia de la Nacion [SCJN], 2016).}
  \accentbar
\end{frame}

\begin{frame}{Seguridad juridica transgredida}
  \begin{columns}[T]
    \begin{column}{0.48\textwidth}
      \begin{block}{Que se transgrede}
        La seguridad juridica se afecta cuando una persona no puede confiar en
        que la autoridad respetara libertad, integridad, legalidad, busqueda,
        investigacion y acceso a la justicia.
      \end{block}
      \begin{block}{Como se observa}
        La falla no esta solo en el hecho inicial, sino en la respuesta
        posterior: verdad incompleta, busqueda insuficiente, coordinacion
        irregular y dano prolongado a familias.
      \end{block}
    \end{column}
    \begin{column}{0.48\textwidth}
      \begin{alertblock}{Por quien}
        La transgresion se atribuye al aparato estatal en sus deberes de
        prevenir, proteger, buscar, investigar, coordinar y reparar. No se
        agota en una persona: revela una falla institucional.
      \end{alertblock}
      \begin{block}{Consecuencia}
        La norma sigue existiendo, pero su eficacia protectora queda en duda.
      \end{block}
    \end{column}
  \end{columns}
  \sourcecite{(Camara de Diputados, 2024, 2026a; CNDH, 2024).}
  \accentbar
\end{frame}

\begin{frame}{Derechos humanos violentados}
  \small
  \begin{tabularx}{\linewidth}{@{}p{0.25\linewidth}Y Y@{}}
    \toprule
    \textbf{Derecho} & \textbf{Afectacion en el caso} & \textbf{Lectura desde universalidad} \\
    \midrule
    Libertad personal & Desaparicion y privacion de control sobre la propia persona. & Ninguna autoridad puede tratar a la persona como objeto disponible. \\
    Integridad y vida & Riesgo extremo, dano fisico, psicologico y social. & La proteccion no depende de origen social, edad, oficio o lugar. \\
    Verdad y justicia & Incertidumbre prolongada y respuesta institucional insuficiente. & Las victimas tienen derecho a saber, exigir y ser escuchadas. \\
    Dignidad humana & Familias y estudiantes son reducidos a dano administrado. & La dignidad opera como fundamento comun de todos los derechos. \\
    \bottomrule
  \end{tabularx}
  \sourcecite{(Camara de Diputados, 2026a; CNDH, 2024; SCJN, 2016).}
  \accentbar
\end{frame}

\begin{frame}{Importancia de los derechos humanos en Mexico}
  \begin{columns}[T]
    \begin{column}{0.31\textwidth}
      \begin{block}{Limite}
        Funcionan como frontera frente al poder: ninguna politica publica,
        investigacion o razon de seguridad puede cancelar la dignidad.
      \end{block}
    \end{column}
    \begin{column}{0.31\textwidth}
      \begin{block}{Metodo}
        Obligan a razonar con persona, dano, deber institucional, reparacion y
        no repeticion; no solo con expediente.
      \end{block}
    \end{column}
    \begin{column}{0.31\textwidth}
      \begin{block}{Confianza}
        Permiten evaluar si el Derecho es una garantia real o solo una promesa
        declarativa frente a la violencia.
      \end{block}
    \end{column}
  \end{columns}
  \vspace{0.28cm}
  \begin{alertblock}{Cierre del eje}
    Como explica Ansuategui Roig (2007), en Mexico los derechos humanos
    importan porque hacen exigible que el Estado responda con verdad, justicia,
    reparacion y control del poder.
  \end{alertblock}
  \sourcecite{(Ansuategui Roig, 2007; SCJN, 2016).}
  \accentbar
\end{frame}

\unadmsection{Propuestas y cierre}

\begin{frame}{Propuestas para prevenir casos similares}
  \begin{columns}[T]
    \begin{column}{0.48\textwidth}
      \begin{block}{1. Busqueda inmediata coordinada}
        Registros interoperables, alertas tempranas y responsables claros desde
        las primeras horas. Fin: eficacia, verdad y proteccion de vida.
      \end{block}
      \begin{block}{2. Investigacion independiente}
        Cadena de custodia, control judicial, sancion por ocultamiento y
        proteccion a testigos. Fin: justicia correctiva y seguridad juridica.
      \end{block}
    \end{column}
    \begin{column}{0.48\textwidth}
      \begin{block}{3. Reparacion integral}
        Atencion psicologica, informacion constante, medidas economicas y
        garantias de no repeticion. Fin: equidad y dignidad.
      \end{block}
      \begin{block}{4. Control democratico}
        Capacitacion, supervision civil, transparencia y protocolos de uso de
        fuerza. Fin: bien comun y limite al poder.
      \end{block}
    \end{column}
  \end{columns}
  \sourcecite{(Camara de Diputados, 2024, 2026b; CNDH, 2024).}
  \accentbar
\end{frame}

\begin{frame}{Postura personal}
  \begin{columns}[T]
    \begin{column}{0.48\textwidth}
      \begin{block}{Posicion}
        El caso Ayotzinapa demuestra la insuficiencia de una lectura puramente
        formal del Derecho. Un sistema juridico no se defiende solo porque
        tenga normas escritas.
      \end{block}
      \begin{block}{Razon}
        Se defiende cuando esas estructuras actuan como limite frente a la
        arbitrariedad y como garantia efectiva para las victimas.
      \end{block}
    \end{column}
    \begin{column}{0.48\textwidth}
      \begin{alertblock}{Consecuencia}
        Si el Estado falla justo donde debe proteger con mayor firmeza, el
        problema deja de ser marginal y se vuelve estructural: afecta la
        legitimidad del orden juridico.
      \end{alertblock}
      \begin{block}{Criterio profesional}
        Como advierten Ronquillo Armas (2018) y Singer (1995), argumentar exige
        encadenar hecho, derecho afectado, norma aplicable, deber incumplido,
        dano y reparacion.
      \end{block}
    \end{column}
  \end{columns}
  \sourcecite{(Ronquillo Armas, 2018; Singer, 1995; Laun, 2021).}
  \accentbar
\end{frame}

\begin{frame}{Conclusion}
  \begin{block}{Hallazgo}
    Ayotzinapa permite comprender que la Filosofia del Derecho no es un plano
    abstracto separado de la realidad: sirve para evaluar si el orden juridico
    conserva sentido cuando enfrenta sus pruebas mas duras.
  \end{block}
  \begin{block}{Sintesis}
    Derecho, moral, justicia, equidad y dignidad no pueden analizarse por
    separado cuando el problema involucra desaparicion de personas, impunidad y
    falla institucional.
  \end{block}
  \begin{alertblock}{Ultima idea}
    El valor de este analisis no esta en repetir que hubo una violacion grave,
    sino en mostrar como esa violacion pone a prueba todo el sistema y obliga a
    tomar posicion frente a el.
  \end{alertblock}
  \sourcecite{(Ruiz Rodriguez, 2009; Rojas-Gonzalez, 2018; Camara de Diputados, 2026a; CNDH, 2024).}
  \accentbar
\end{frame}

\begin{frame}{Referencias bibliograficas (1/2)}
  \scriptsize
  \begin{itemize}
    \setlength\itemsep{0.10cm}
    \item Ansuategui Roig, F. J. (2007). \textit{De los derechos y el Estado de derecho: aportaciones a una teoria juridica de los derechos}. Universidad Externado de Colombia.
    \item Camara de Diputados del H. Congreso de la Union. (2024). \textit{Ley General de Victimas}. Diario Oficial de la Federacion. \url{https://www.diputados.gob.mx/LeyesBiblio/pdf/LGV.pdf}
    \item Camara de Diputados del H. Congreso de la Union. (2026a). \textit{Constitucion Politica de los Estados Unidos Mexicanos}. Texto vigente; ultima reforma DOF 03-03-2026. \url{https://www.diputados.gob.mx/LeyesBiblio/ref/cpeum.htm}
    \item Camara de Diputados del H. Congreso de la Union. (2026b). \textit{Ley General en Materia de Desaparicion Forzada de Personas, Desaparicion Cometida por Particulares y del Sistema Nacional de Busqueda de Personas}. Diario Oficial de la Federacion. \url{https://www.diputados.gob.mx/LeyesBiblio/pdf/LGMDFP.pdf}
    \item Comision Nacional de los Derechos Humanos. (2024). \textit{Recomendacion No. VG 2024 sobre el caso Ayotzinapa}. \url{https://www.cndh.org.mx/documento/recomendacion-no-vg-2024-sobre-el-caso-ayotzinapa}
    \item Kelsen, H. (1982). \textit{Teoria pura del derecho} (R. J. Vernengo, Trad.). Universidad Nacional Autonoma de Mexico.
  \end{itemize}
  \accentbar
\end{frame}

\begin{frame}{Referencias bibliograficas (2/2)}
  \scriptsize
  \begin{itemize}
    \setlength\itemsep{0.10cm}
    \item Laun, R. (2021). \textit{Derecho y moral}. Ediciones Olejnik. \url{https://doi.org/10.5771/9789563929119}
    \item Rojas-Gonzalez, G. (2018). \textit{Filosofia del Derecho}. Universidad Catolica de Colombia.
    \item Ronquillo Armas, L. A. (2018). \textit{Etica general y profesional} (2.a ed.). Editorial Mar y Trinchera.
    \item Ruiz Rodriguez, V. (2009). \textit{Filosofia del derecho}. Instituto Electoral del Estado de Mexico.
    \item Singer, P. (Ed.). (1995). \textit{Compendio de etica}. Alianza Editorial.
    \item Suprema Corte de Justicia de la Nacion, Primera Sala. (2016). \textit{Dignidad humana. Constituye una norma juridica que consagra un derecho fundamental a favor de las personas. Jurisprudencia 1a./J. 37/2016 (10a.)}. Registro digital: 2012363. \url{https://sjf2.scjn.gob.mx/detalle/tesis/2012363}
    % \item Universidad Abierta y a Distancia de Mexico. (2026). \textit{Planificacion de actividades S8 - Filosofia del Derecho} [Planeacion de actividad].
  \end{itemize}
  \accentbar
\end{frame}
% \begin{frame}{Fuentes de consulta}
%   \scriptsize
%   \begin{tabularx}{\linewidth}{@{}p{0.30\linewidth}Y@{}}
%     \toprule
%     \textbf{Fuente} & \textbf{Uso en la presentacion} \\
%     \midrule
%     Camara de Diputados. Constitucion Politica de los Estados Unidos Mexicanos. & Derechos humanos, seguridad juridica, acceso a justicia y derechos de victimas. \\
%     Camara de Diputados. Ley General de Victimas. & Ayuda, asistencia, verdad, reparacion integral y garantias de no repeticion. \\
%     Camara de Diputados. Ley General en Materia de Desaparicion Forzada de Personas. & Busqueda, investigacion y coordinacion institucional. \\
%     CNDH. Recomendacion institucional sobre el caso Ayotzinapa. & Hechos, relevancia publica y afectacion a victimas. \\
%     SCJN, Primera Sala. Jurisprudencia 1a./J. 37/2016. & Dignidad humana como norma juridica fundamental. \\
%     Ruiz Rodriguez; Rojas Gonzalez; Laun; Ronquillo Armas; Singer. & Fundamento filosofico: derecho, moral, justicia, legitimidad y responsabilidad etica. \\
%     \bottomrule
%   \end{tabularx}
%   \vspace{0.15cm}
%   {\color{unadmMuted}Formato APA 7 completo desarrollado en el reporte base y en \texttt{filosofia-del-derecho.bib}.}
%   \accentbar
% \end{frame}

% \begin{frame}{Control para grabar el video}
%   \begin{columns}[T]
%     \begin{column}{0.48\textwidth}
%       \begin{block}{Requisitos de contenido}
%         \begin{itemize}
%           \item Caratula, introduccion, desarrollo, propuestas, conclusion y fuentes.
%           \item Un caso controversial de Mexico: Ayotzinapa.
%           \item Constitucion mexicana y al menos dos leyes relacionadas.
%           \item Lenguaje academico, sin lectura completa.
%         \end{itemize}
%       \end{block}
%     \end{column}
%     \begin{column}{0.48\textwidth}
%       \begin{alertblock}{Requisitos tecnicos}
%         Duracion de 4 a 8 minutos, buena iluminacion, audio claro, camara
%         encendida, rostro visible, vestimenta formal o semiformal y sin musica
%         de fondo.
%       \end{alertblock}
%       \begin{block}{Entrega}
%         Video MP4 o enlace; si se entrega enlace, agregar PDF con portada, enlace
%         y 3 a 5 fotos del video.
%       \end{block}
%     \end{column}
%   \end{columns}
%   \accentbar
% \end{frame}

\end{document}
